\documentclass[10pt, a4paper, twocolumn]{article}

\usepackage[utf8]{inputenc}
\usepackage[T1]{fontenc}
\usepackage[spanish]{babel}
\usepackage{amsmath, amssymb, amsthm}
\usepackage{geometry}
\usepackage{booktabs}
\usepackage{graphicx}
\usepackage{float}
\usepackage{hyperref}
\usepackage{parskip}

\geometry{margin=2cm}

% Recuadro para marcar donde va cada grafico. Ocupa el ancho completo de la
% hoja (las dos columnas), porque todas las figuras tienen dos columnas de
% subgraficos: la metrica (retorno o pasos) de un lado y la loss del otro.
% Son tres en total, todas generadas en el notebook:
%   1. la figura resumen apilada de los 6 entornos (6 filas x 2 columnas),
%      que sale de plot_resumen_entornos;
%   2. y 3. las comparaciones con y sin baseline de CartPole y Acrobot,
%      que salen de plot_baseline_comparison.
%
% OJO: no se fuerza un salto de pagina despues de cada una (con 3 imagenes mas
% todo el texto que pide la consigna, forzarlo garantiza pasarse de 3
% paginas). Quedan flotando arriba de la pagina que les toque. Si en Overleaf
% dos terminan cayendo en la misma hoja y no queda bien, se puede meter un
% \clearpage puntual entre esas dos nada mas.
% El primer argumento es opcional y es la etiqueta para poder referenciarla con
% \ref. Tiene que ir DENTRO del figure* y despues del \caption: si se pone un
% \label suelto afuera, \ref devuelve el numero de la seccion y no el de la
% figura.
\newcommand{\imgancha}[4][]{%
\begin{figure*}[t]
\centering
\fbox{\begin{minipage}[c][#3][c]{0.94\textwidth}
\centering\footnotesize #2
\end{minipage}}
\caption{\footnotesize #4}
\ifx\relax#1\relax\else\label{#1}\fi
\end{figure*}
}

\title{
    \textbf{Trabajo Práctico 1} \\
    \large Gradiente de Política \\
    \normalsize I-404 -- Aprendizaje reforzado \\
    Universidad de San Andrés
}
\author{Carlos Federico Trejo}  % <-- completar
\date{Septiembre 2026}

\begin{document}

\maketitle

% ========================================.

\section{Introducción}

En este trabajo se implementó el algoritmo REINFORCE (gradiente de política) y se lo usó para entrenar agentes en seis entornos distintos: cuatro chicos armados especialmente para poder debuggear el algoritmo con tranquilidad y dos que ya vienen con Gymnasium (CartPole-v1 y Acrobot-v1), más exigentes. Después se extendió el algoritmo agregándole un baseline (una red que estima cuánto retorno se espera desde cada estado) para tratar de reducir la varianza del entrenamiento. En este informe se cuenta cómo se armó, qué resultados dio, y qué cosas se aprendieron en el camino que no se esperaban de antemano.

\section{Implementación de REINFORCE}

La política se representa con una red neuronal chica (dos capas ocultas de 20 neuronas con activación ELU) que recibe el estado y devuelve un puntaje (logit) por cada acción posible. Para elegir la acción no se usa el puntaje más alto: se arma una distribución categórica a partir de esos puntajes (con \texttt{torch.distributions.Categorical}) y se muestrea de ahí, porque la política tiene que ser estocástica para poder explorar durante el entrenamiento.

Para entrenar, se junta un batch de varias trayectorias completas con la política actual, se calcula el retorno descontado de cada paso (con $\gamma = 0.99$ en todos los entornos) y esos retornos se usan para armar el gradiente que se optimiza con Adam. Antes de actualizar la política se normalizan los retornos del batch completo (media 0, desvío 1), como recomienda el enunciado, para que el entrenamiento no sea tan ruidoso.

\section{Para qué sirve cada entorno de prueba}

Antes de entrenar en CartPole o Acrobot, se armaron cuatro entornos chicos pensados para ir probando el algoritmo de a poco. La idea es que si algo falla en uno de estos entornos simples, es mucho más fácil darse cuenta de qué parte del código está mal.

\textbf{TwoAZeroObsOneStep.} Es el caso más simple posible: un solo estado (siempre el mismo) y un solo paso por trayectoria. Sirve para probar lo más básico de REINFORCE, que el update realmente aumente la probabilidad de la acción que da mejor recompensa. Si esto no funciona acá, no tiene sentido seguir probando en algo más difícil.

\textbf{TwoARandomObsOneStep.} Acá la acción correcta depende de en qué estado arrancó la trayectoria (S1 o S2). Sirve para confirmar que la red neuronal realmente está usando la observación para decidir qué hacer, y no aprendiendo una sola acción fija de memoria. Una política que ignore el estado no puede resolver este entorno, porque en promedio le va a ir mal la mitad de las veces.

\textbf{LineWorldEasyEnv.} Acá la recompensa aparece recién al final del pasillo, no en cada paso. Este entorno sirve para probar que el algoritmo propaga bien la recompensa hacia atrás en el tiempo, es decir que el cálculo de los retornos descontados (con el factor $\gamma$) está bien hecho. Si el descuento estuviera mal calculado, el agente tendría problemas para darse cuenta de que moverse a la derecha desde el principio del pasillo también conviene aunque el premio esté varios pasos más adelante.

\textbf{LineWorldMirrorEnv.} En este entorno ninguna política que siempre haga lo mismo puede resolverlo, porque justo en el medio del pasillo las acciones se invierten. Sirve para confirmar que el agente aprende una política que realmente depende del estado a lo largo de varios pasos, combinando lo que se prueba en TwoARandomObsOneStep (usar el estado) con lo que se prueba en LineWorldEasyEnv (varios pasos seguidos).

\section{Resultados de entrenamiento}
\label{sec:resultados}

Para los seis entornos se entrenó hasta que el agente resolviera el problema, graficando la curva de recompensa (o de pasos por trayectoria, en los dos pasillos, donde lo que importa es llegar en menos pasos) y la curva de loss, más las líneas de referencia de un agente aleatorio y de los agentes triviales que pide la consigna. En la figura~\ref{fig:resumen} están los seis juntos, una fila por entorno: a la izquierda la métrica de desempeño con sus líneas de referencia, y a la derecha la loss. Se muestran apilados y no como seis figuras sueltas simplemente porque el informe tiene un límite de tres páginas; en el notebook cada entorno tiene además su propio gráfico por separado, más grande.

Lo que se ve es que en los cuatro entornos chicos el agente supera claramente al aleatorio: en los dos de un solo paso el retorno se va a 1, que es el máximo posible, y en los dos pasillos la cantidad de pasos por trayectoria cae hasta el mínimo. Contra los agentes triviales la comparación cambia según el entorno. En LineWorldEasyEnv la política aprendida termina empatando con la acción constante ``siempre a la derecha'' (5 pasos por trayectoria), pero eso es lo esperable: en ese entorno ir siempre a la derecha ya es óptimo, y lo que se está comprobando es que el agente lo descubra solo. En TwoARandomObsOneStep las dos acciones constantes quedan en $0{,}11$ y $-0{,}11$, mientras que el agente llega a 1, que es lo que tenía que pasar en un entorno donde la acción correcta depende del estado. Y en LineWorldMirrorEnv directamente no hay ninguna línea de referencia trivial en el gráfico: ninguna de las dos acciones constantes llega nunca a la meta, así que no hay con qué comparar más allá del agente aleatorio. En CartPole-v1 el retorno llega al techo de 500 pasos del entorno, y en Acrobot-v1 pasa de los $-500$ del agente aleatorio (que nunca llega a levantar el péndulo y siempre se le termina el tiempo) a alrededor de $-100$.

Un detalle que al principio confundió es que las curvas de loss no bajan de forma prolija como uno esperaría de una red entrenada de la manera clásica, y en CartPole y Acrobot ni siquiera tienden a cero. No es un error: en REINFORCE la loss no es un error que haya que minimizar, es solo un artificio para que el optimizador haga ascenso de gradiente, y su valor depende de los retornos del batch que se juntó en esa iteración, que van cambiando a medida que la política mejora. Por eso el gráfico que hay que mirar para saber si el agente está aprendiendo es el de la izquierda, no el de la derecha.

\imgancha[fig:resumen]{Gráfico: figura resumen 6x2 -- los seis entornos apilados, métrica de desempeño (izquierda) y loss (derecha)}{17.5cm}{Los seis entornos entrenados con REINFORCE (Parte 3). Una fila por entorno: a la izquierda retorno por trayectoria (o pasos por trayectoria en los dos \textit{LineWorld}) junto con las líneas de referencia del agente aleatorio y de los agentes triviales, a la derecha la loss. En todos los casos la línea clara es el valor trayectoria a trayectoria y la oscura su media móvil.}

Para CartPole-v1 y Acrobot-v1 se muestran además, más abajo, los gráficos que comparan la versión con y sin baseline (Parte 4), donde la curva de la Parte 3 aparece superpuesta como una de las dos líneas.

\imgancha{Gráfico: CartPole-v1 -- retorno y loss por trayectoria, con y sin baseline}{5.5cm}{CartPole-v1: comparación con baseline (Parte 4) y sin baseline (Parte 3).}

\imgancha{Gráfico: Acrobot-v1 -- retorno y loss por trayectoria, con y sin baseline}{5.5cm}{Acrobot-v1: comparación con baseline (Parte 4) y sin baseline (Parte 3).}

\section{El entorno que más costó}

Los seis entornos se entrenaron con el mismo presupuesto (150 iteraciones de un batch de 16 trayectorias, o sea 2400 trayectorias cada uno), justamente para que se puedan comparar entre sí. Con presupuestos distintos, preguntarse cuál costó más termina midiendo cuánto se decidió entrenar cada uno y no qué tan difícil es.

Sobre esas curvas se miden dos cosas distintas, que no dan lo mismo. \textbf{Aprender} es en qué trayectoria la curva suavizada recorre por primera vez el 90\% del camino entre donde arrancó y su meseta. \textbf{Estabilizarse} es a partir de qué trayectoria ya no se sale más de una banda del 5\% alrededor de esa meseta.

\begin{table}[H]
\centering
\footnotesize
\begin{tabular}{lrr}
\toprule
Entorno & Aprende en & Se estabiliza en \\
\midrule
LineWorldEasyEnv & 244 & 265 \\
TwoAZeroObsOneStep & 279 & 284 \\
TwoARandomObsOneStep & 334 & 355 \\
Acrobot-v1 & 399 & \textbf{1359} \\
CartPole-v1 & 456 & 478 \\
LineWorldMirrorEnv & \textbf{467} & 527 \\
\bottomrule
\end{tabular}
\end{table}

El que más tarda en aprender es \textbf{LineWorldMirrorEnv}. Tiene sentido, porque es el único entorno donde no hay ningún atajo: en LineWorldEasyEnv, por ejemplo, con solo aprender a ir siempre para la derecha ya se resuelve todo. En LineWorldMirrorEnv eso no funciona, porque en el estado del medio hay que hacer exactamente lo contrario de lo que conviene en los otros. El agente tiene que aprender una política que distinga entre estados, no una preferencia general por una acción, y eso le lleva más trayectorias.

El caso de \textbf{Acrobot-v1} es distinto y vale la pena separarlo. Aprende rápido (llega al 90\% de su meseta en la trayectoria 399, antes que CartPole y que los dos \textit{LineWorld}), pero tarda más del triple en estabilizarse que cualquier otro. Esa diferencia entre las dos columnas no es dificultad: es varianza. La curva sigue teniendo bajones mucho después de haber alcanzado un buen retorno promedio, y reducir esa varianza es exactamente lo que busca el baseline de la Parte 4.

Si en cambio se mira la cantidad total de pasos de simulación en vez de trayectorias, el que más costó fue CartPole-v1, porque a medida que el agente mejora las trayectorias se hacen mucho más largas (hasta 500 pasos cada una), entonces cada iteración termina siendo mucho más pesada de correr.

\section{Qué fue lo más demandante en tiempo}

Lo que más tiempo demandó no fue programar el algoritmo en sí, sino lograr que las corridas largas de entrenamiento no se cuelguen al ejecutar el notebook completo en Windows por problemas de comunicación con el kernel de Jupyter; en retrospectiva, se podría haber evitado corriendo los entrenamientos como scripts de Python sueltos y guardando los resultados aparte, en vez de depender de que todo el notebook se ejecute de punta a punta sin cortes.

\section{Comparación con baseline (Parte 4)}

El baseline se implementó como una segunda red, más chica que la de la política (una sola capa oculta), que recibe el estado y devuelve un único número: una estimación de $V(s)$, o sea cuánto retorno se espera desde ahí en adelante. Con eso se arma la ventaja $A(t) = R(t) - V(s_t)$, que es lo que se usa en vez del retorno directo para actualizar la política, y la red de valor se entrena aparte, por regresión, para que $V(s_t)$ se acerque al retorno real $R(t)$.

Las curvas de CartPole-v1 y Acrobot-v1 de la sección~\ref{sec:resultados} ya muestran la comparación entre la versión con baseline (red de valor $V(s)$, Parte 4) y la versión sin baseline (Parte 2), superpuestas. La mejora que se ve ahí es más chica de lo que se esperaba en un principio, y se cree que es porque la versión sin baseline ya normaliza los retornos del batch (le resta la media y divide por el desvío antes de actualizar la política), y restar la media es, en el fondo, un caso particular de baseline: uno que no depende del estado, en vez de uno que sí depende como $V(s)$. Por eso la diferencia entre las dos versiones es más chica que si se comparara contra una versión de REINFORCE que no reste ni normalice nada.

\section{Conclusión}

Lo que más queda de este trabajo es que el algoritmo en sí no es tan largo de escribir, pero hace falta mucho cuidado para confiar en que está bien implementado: los entornos de prueba de la Parte 1 fueron clave para eso, porque permitieron confirmar que cada parte del algoritmo (el signo del gradiente, el uso de la observación, la propagación del retorno en el tiempo) funcionaba bien por separado, en entornos donde un error se nota enseguida, antes de meterlo todo junto en CartPole o Acrobot. También sorprendió que agregar el baseline no mejorara tan claramente como se esperaba según la consigna, y se terminó entendiendo que era porque la normalización de retornos de la Parte 2 ya cumplía en parte ese mismo rol.

\section{Videos}

Los videos de evaluación de \texttt{CartPole-v1} y \texttt{Acrobot-v1} (una trayectoria completa, política final de forma determinista) están incluidos en el .zip entregado, en la carpeta \texttt{videos/}. % <-- opcional: reemplazar por link de Drive si se prefiere

\end{document}
